\documentclass[10pt,a4paper]{article}

\usepackage[left=1.45cm,right=1.45cm,top=1.2cm,bottom=1.2cm]{geometry}
\usepackage{enumitem}
\usepackage{hyperref}
\usepackage{xcolor}
\usepackage{titlesec}
\usepackage{fontspec}

\definecolor{darkblue}{HTML}{0F172A}
\definecolor{accentblue}{HTML}{2563EB}
\hypersetup{colorlinks=true,urlcolor=accentblue,linkcolor=darkblue}
\setlength{\parindent}{0pt}
\setlength{\parskip}{0pt}
\renewcommand{\familydefault}{\sfdefault}
\linespread{1.02}
\raggedbottom
\titleformat{\section}{\normalsize\bfseries\color{darkblue}}{}{0pt}{}[\titlerule]
\titlespacing*{\section}{0pt}{3pt}{2pt}
\setlist[itemize]{leftmargin=1.15em,itemsep=1pt,topsep=2pt,parsep=0pt,partopsep=0pt}

\begin{document}

\begin{center}
    {\LARGE \textbf{ANDRÉ NICOLAS SILVA}}\\[-2pt]
    {\normalsize \textbf{Analista de Operações Jr. | Processos, Automação e Tecnologia}}\\[2pt]
    \footnotesize Maceió, AL \quad • \quad
    \href{mailto:devandrenicolas@gmail.com}{devandrenicolas@gmail.com}\\
    \href{https://portfolio-andrenicolas.vercel.app}{Portfólio}
    \quad • \quad \href{https://github.com/devAndreNicolas}{GitHub}
    \quad • \quad \href{https://www.linkedin.com/in/devandrenicolas}{LinkedIn}
\end{center}

\section*{RESUMO PROFISSIONAL}
Profissional de tecnologia com formação em Sistemas para Internet e experiência em acompanhar funcionalidades de software desde a análise de necessidades e regras até testes, depuração, entrega e suporte. Atua em uma plataforma SaaS usada por mais de 10 empresas, com prática em organização de entregas, investigação de falhas, documentação técnica e colaboração com produto e engenharia. Familiaridade prática com Airtable, Make, n8n, Notion e Linear. Traz competências de organização, suporte técnico e comunicação de bloqueios, além de base técnica para aprender processos e ferramentas operacionais.

\section*{COMPETÊNCIAS}
\begin{itemize}
    \item \textbf{Operações e processos:} organização de tarefas e registros, acompanhamento de pendências, documentação técnica, comunicação de bloqueios, atendimento e escalonamento como competências declaradas.
    \item \textbf{Ferramentas:} Airtable, Make, n8n, Notion e Linear; ferramentas de outbound (familiaridade prática); Excel e Microsoft Office; Power BI para análise e relatórios.
    \item \textbf{Qualidade e suporte:} acompanhamento de testes, registro de resultados, troubleshooting de aplicações, análise de erros, suporte à produção e SQL para investigação de dados.
    \item \textbf{Tecnologia e automação:} integrações, automações entre sistemas, APIs, SQL/PostgreSQL, workflows de software e especificações estruturadas (SDD).
    \item \textbf{Idiomas:} Português nativo; Inglês intermediário/avançado, em aprimoramento contínuo.
\end{itemize}

\section*{EXPERIÊNCIA}
\textbf{MSPA — Frontend Engineer (contribuições full stack)} \hfill \textit{Abr 2025 -- Presente}\\
\textit{Plataforma SaaS para privacidade, conformidade com a LGPD e auditoria digital}
\begin{itemize}
    \item Acompanha funcionalidades do Compass, Trust Center e outras áreas da plataforma, conectando necessidades de produto e regras de negócio a escopo técnico, implementação, validação e suporte; os produtos são usados por mais de 10 empresas.
    \item Investiga problemas em fluxos entre interface, APIs, serviços e dados, usando depuração e SQL/PostgreSQL para entender comportamentos e apoiar a resolução.
    \item Participa do ciclo de entrega, incluindo testes automatizados, documentação técnica, revisão de mudanças, depuração em produção e integração entre camadas do sistema.
    \item Contribui com integrações e automações de software conforme as necessidades dos produtos, mantendo atenção às regras de negócio e ao comportamento esperado.
    \item Reorganizou o fluxo de trabalho em uma fase crítica de um projeto com seis pessoas, coordenando diretamente quatro integrantes e apoiando a clareza de escopo e a organização das entregas.
\end{itemize}

\section*{PROJETOS RELEVANTES}
\textbf{Marca Racha — Fluxos e regras de operação} \hfill \textit{Next.js, TypeScript, Supabase, PostgreSQL, Vitest}\\
Desenvolveu fluxos para organização de partidas e solicitações de entrada. Testes comportamentais verificam permissões, visibilidade, notificações e transições de estado, ajudando a registrar e conferir resultados esperados.

\vspace{2pt}
\textbf{Compass e Trust Center — Plataforma SaaS} \hfill \textit{Angular, TypeScript, Lit, Web Components}\\
Participou da análise de fluxos de produto, entrevistas e validações de UX, além de instrumentação e análise de eventos e funis no PostHog para apoiar decisões sobre a experiência do produto.

\section*{FORMAÇÃO}
\textbf{Tecnólogo em Sistemas para Internet} — UNCISAL \hfill \textit{Concluído em 2026}

\end{document}
